% Backing chapter for the claim "who can do what with AI" (slug who-gains).
% Author's note (verification and reproduction): all searches are in the
% scholar session `who-gains` (27,431 records after the repair queries);
% the inherited session `who-gains-fulltext` (OpenAlex full-text
% matching, capped or failed) was not used. Query strings, the screening
% batches and their validated outputs, and the provenance of every cited
% source are in litteratursokning/who-gains-report.md,
% litteratursokning/who-gains-screen/ and litteratursokning/who-gains.bib.
% Status: screened by AI agents; no human sample review yet.

\chapter{Who gains more from generative-AI assistance at work: less-skilled or novice workers, or more-skilled and expert workers?}
\label{who-gains:chapter}

\emph{Status: screened by an AI agent; awaiting human review.}

The lecture says that what a person gets out of a generative-AI
assistant depends on what that person can do without it. Its working
thesis was that AI is a \enquote{skill amplifier}: it complements
competent people and amplifies incompetence too. If that is true, the
more skilled should gain more from the assistance than the less skilled,
and the least skilled might even be harmed. The opposite hypothesis is
that AI levels the field: novices gain most, because the model supplies
what they lack. This chapter tests the thesis in both directions. It
asks which of the two the empirical studies of 2022--2026 support, for
which tasks, and how strongly.

\section{Method}
\label{who-gains:method}

\paragraph{Scope.}
The review covered empirical studies (experiments, field and
quasi-experiments, observational studies with outcome data, systematic
reviews and meta-analyses) of assistance based on large language models
in knowledge work, programming included. Studies had to be published
2022--2026 and report how the effect differs by the worker's skill,
experience, tenure or baseline ability. Studies of what moderates that
difference were also included: the kind of task, whether the task lies
inside what the model does well, the ability to judge or verify the
output, and measured versus self-reported gains. Studies of students
count when they measure performance on a work-like task split by prior
skill.

\paragraph{Databases and their health.}
All searches ran on 27 November 2026. They went to Semantic Scholar
(S2), OpenAlex (OA, not open access), Scopus, Web of Science (WoS),
IEEE Xplore (IEEE) and DBLP.
\begin{itemize}
\item DBLP answered every query with a page that stops automated clients
  instead of with data, so it could not be searched.
\item IEEE spent its daily quota before the last three queries.
\item The searches ran without the university network, so Scopus
  returned no abstracts. Full texts were fetched afterwards through the
  university network where the publisher allowed it; two publishers
  answered with a bot check instead, and those sources stay at abstract
  level.
\item No search had an upper limit on the number of hits.
\item OpenAlex matched title and abstract only. Semantic Scholar answered
  Boolean queries through its bulk search.
\end{itemize}

\paragraph{Queries.}
Seven queries share three facets: generative AI, skill level, and
knowledge work or workers. All but one add a fourth facet for outcomes
(productivity, performance, quality). Each then adds its own facet, as
\cref{who-gains:queries} shows. The supporting queries (S) are phrased
to surface amplification or harm to the less skilled. The counter-queries
(C) are phrased to surface levelling, heterogeneity in general,
moderators, self-reported gains and reviews. Every concept has several
wordings (for example \enquote{equalise}, \enquote{level},
\enquote{skill gap}, \enquote{narrow}, \enquote{compress}, \enquote{low
performers}). Every query ran with the same keywords on every database:
WoS received them as a topic search, Scopus as a title-abstract-keyword
search, the others as Boolean text.

\paragraph{Known studies and the check by meaning.}
Six studies were named before the searches as ones that must come back.
They were field or randomised studies of customer support
\autocite{Brynjolfsson2025Generative}, professional writing
\autocite{Noy2023Experimental}, consulting \autocite{DellAcqua2026Navigating},
software development \autocite{Cui2026Effects}, entrepreneurs
\autocite{Otis2026Uneven} and experienced open-source developers
\autocite{Becker2025Measuring}. The keyword queries were then checked
against a semantic search, with four questions of 50 works each. The
questions were on levelling, amplification, expert slowdown and the task
frontier. The check showed two gaps. Randomised studies that speak of
\enquote{AI tools}, \enquote{AI advice} or \enquote{AI mentors} were
missing, and so were experiments that describe skill as
\enquote{ability} or \enquote{junior} and \enquote{senior}. Two repair
queries by concept closed both gaps (C6 and C7 in
\cref{who-gains:queries}). After them, all six named studies were in the
session. Four relevant works that the semantic search returned remained
outside every keyword answer, because no database held an abstract for
them. They are listed under future work. One further experiment came
back only through the semantic search, because its abstract names no
skill term \autocite{Peng2023Impact}.

\begin{table}[tbp]
  \begin{sidecaption}{Queries and hits per database, 27 November 2026.
    S = supporting query (phrased towards amplification), C =
    counter-query (phrased towards levelling or a moderator).
    Databases: Semantic Scholar (S2), OpenAlex (OA; not open access),
    Scopus, Web of Science (WoS), IEEE Xplore (IEEE).  \enquote{N/M}
    means N records returned of M reported.  \enquote{--} means that the
    search failed: IEEE had spent its daily quota.  DBLP failed every
    query (it returned a page that stops automated clients) and is not
    listed.  \enquote{inc.}\ marks an answer that stopped early (S2 rate
    limit, Scopus paging error); its count is a lower bound.  The
    searches had no upper limit.}[who-gains:queries]
    \centering\footnotesize
    \begin{tabular}{@{}lp{0.34\linewidth}rrrrr@{}}
      \toprule
      & Own facet & S2 & OA & Scopus & WoS & IEEE \\
      \midrule
      S1 & amplify, complement, widen, skill-biased, high performers,
           returns to skill & 614 & 1,473 & 352 & 391 & 125 \\
      S2 & over-reliance, automation bias, slower, worse, harm,
           deskilling, jagged frontier (no outcome facet)
           & 1,060 & 2,194 & 580 & 613 & 173 \\
      C1 & equalise, level, skill gap, narrow, compress, low performers,
           less experienced & 2,000 & 4,118 & 203 & 1,327 & 672 \\
      C2 & heterogeneous, differential, who benefits, largest gains
           & 263 & 606 & 216 & 201 & 71 \\
      C3 & verification, evaluate, judge, task complexity, task type,
           moderator & 3,896 & 7,054 & 1,701 & 2,235 & 945 \\
      C4 & perceived, self-reported, belief, expectations, forecast
           & 1,445 & 3,183 & 1,003 & 1,232 & 403 \\
      C5 & meta-analysis, systematic review, scoping review
           & 378 & 840 & 273 & 220 & -- \\
      C6 & repair: broader AI wording (AI tools, AI advice, AI mentor)
           $\times$ skill $\times$ randomised, field or natural
           experiment, difference-in-differences $\times$ outcome
           & 94 & 196 & 98 & 108 & -- \\
      C7 & repair: broader AI wording $\times$ ability, junior, senior,
           performance gap, expertise $\times$ experiment or trial
           $\times$ outcome
           & 3,995 inc. & 9,477 & 4,925 inc. & 4,035 & -- \\
      \bottomrule
    \end{tabular}
  \end{sidecaption}
\end{table}

\paragraph{Screening.}
The queries returned 27,431 distinct records after deduplication. Copies
of a work with the same identifier were merged into one record, and a
preprint was merged with its published version where the match was
unambiguous. Screening went in tiers:
\begin{itemize}
\item \emph{Tier 1} held the first 100 records of every answer ordered by
  relevance. To these were added every work of the semantic check that
  was in the session, and the named studies. That makes 1,694 records.
\item AI agents screened tier 1 in batches of at most 100. For each
  record they gave a verdict, one of four categories, a labelled
  exclusion criterion for every discard, and a reason quoting the
  abstract. Every batch passed an automatic check for defaulted rows and
  for quotes not found in the record.
\item From the 13,340 records outside tier 1 before the repair queries,
  a seeded random sample of 60 was screened the same way.
\item The remaining 25,577 records were counted but not screened or
  classified. Only a single trial batch of 100 of them was
  classified by a language model from the stated research question
  (99 discarded, 1 kept); the hit list includes those 100 decisions.
\end{itemize}
The screening therefore covers tier 1 in full, the sample of 60 and the
trial batch, not the whole set of records. No human has yet reviewed a
sample of these decisions, so their error rate is unknown. Each source
cited below was opened by the agent. Where the full text could not be
obtained, this is said where the source is used. Every quotation is
verbatim.

\section{Results}
\label{who-gains:results}

\paragraph{The screen.}
Of the 1,694 tier-1 records, 103 bore on the question and were kept.
Of these, 23 point towards amplification and 80 towards levelling or a
moderator. Some works appear as several records (a preprint, a working
paper and a journal version), so the hit list, which counts works,
shows fewer rows. Another 74 bore on the question but were excluded: for
example opinion pieces without data, records without abstracts, and
non-generative AI such as image detectors in endoscopy. Further, 1,334
records were in the same field but did not compare skill levels, and 183
were off topic. Among the kept records, levelling was the most common
theme (56), then amplification (17), task-dependence (9), verification
(7), syntheses (7), no difference (5) and perceived versus measured
gains (2). The random sample of the 13,340 records outside tier 1 before the
repair queries held no record that bore on the question: 0 of 60. The
exact binomial 95\,\% upper bound on the share of those records that
bear on the question is therefore 6.0\,\%, that is, at most about 800
of them. The repair queries added about 12,900 records after the sample
was drawn, so the bound does not cover them. In all, 25,577 records were
counted but neither screened nor classified. The sample is the only
statement this chapter makes about them.

\paragraph{Levelling on well-defined tasks.}
The strongest and most consistent evidence is that the less skilled gain
most on tasks the model handles well.
\begin{itemize}
\item In customer support, \textcite{Brynjolfsson2025Generative} followed
  5,172 agents through a staggered roll-out. \enquote{Less experienced
  and lower-skilled workers improve both the speed and quality of their
  output, while the most experienced and highest-skilled workers see
  small gains in speed and small declines in quality}.
\item Two later customer-service field experiments agree. At Alibaba,
  \enquote{Lower-performing agents benefit the most}
  \autocite{Ni2026Generative}. At a delivery firm, the benefits were most
  pronounced for less-experienced agents \autocite{Zhang2026Engaging}.
\item In professional writing, a preregistered experiment with
  college-educated professionals found that ChatGPT \enquote{compresses
  the productivity distribution by benefiting low-ability workers more}:
  lower-ability workers gained quality and time, while higher-ability
  workers kept their quality and became faster
  (\cite{Noy2023Experimental}, working-paper version; the published
  abstract confirms that \enquote{Inequality between workers
  decreased}).
\item In software development, three company field experiments with
  4,867 developers found that \enquote{less experienced developers had
  higher adoption rates and greater productivity gains}
  \autocite{Cui2026Effects}. Only the abstract could be read: the
  publisher's full text was behind a bot check. An earlier
  controlled experiment found that developers with less programming
  experience benefited most \autocite{Peng2023Impact}.
\item Among 758 consultants on tasks inside the model's competence,
  those below the average baseline improved by 43\,\% and those above by
  17\,\% (\cite{DellAcqua2026Navigating}, working-paper version).
\item The pattern recurs in tutoring, where students of lower-rated
  tutors gained most \autocite{Wang2024Tutor}.
\item It also holds in a controlled experiment in which \enquote{AI
  improves performance more for people with low baseline ability}
  \autocite{Caplin2026ABCs}.
\item In design work, novices kept benefiting during execution while
  experts lost efficiency \autocite{Hou2025DoubleEdged}.
\item In advertising copy, feedback from a language model brought
  non-experts close to experts \autocite{Chen2024Large}.
\end{itemize}
Two syntheses agree. A review of experiments and field data concludes
that \enquote{Novice workers tend to benefit more from LLMs in simple
tasks} \autocite{DelRioChanona2025AI}. A systematic review of 157
studies reports gains \enquote{typically largest for less experienced
workers} \autocite{Khandare2026Short}; it is a single-author preprint.

\paragraph{The gap narrows but stays.}
Levelling is relative, not absolute. In a randomised experiment with
1,174 adults, AI increased performance for all participants,
\enquote{with substantially larger gains for lower-education
individuals}. Even so, \enquote{higher-education participants use the
tool more effectively across several margins}, and once the assistant
was removed \enquote{a sizable education gap re-emerges}
\autocite{Cruces2026Does}. A small classroom study likewise found that
stronger participants keep an absolute advantage while the gap is
compressed \autocite{Bierbrauer2026Is}.

A three-month experiment with 133 patent lawyers separates the two
effects. Quality while using AI rose \enquote{with larger gains among
junior lawyers}. The unassisted advantage three months later, however,
\enquote{was concentrated entirely among senior lawyers}
\autocite{Autor2026Does}. Assisted performance levels; durable skill
does not.

\paragraph{Amplification on open-ended and unverifiable tasks.}
The evidence for amplification comes from tasks where the user must
choose or check what the model offers.
\begin{itemize}
\item In a field experiment, Kenyan entrepreneurs were given a GPT-4
  business adviser. \enquote{low performers did nearly 10\% worse because
  of the AI assistant, whereas high performers may have benefited by over
  15\%} \autocite{Otis2026Uneven}. The published version attributes this
  to which advice the entrepreneurs chose to implement.
\item In a laboratory experiment, ChatGPT lowered mathematics scores
  \enquote{mainly because the low-ability subjects couldn't discern the
  hallucinated answers with the correct ones} \autocite{Kim2024Helping}.
\item On a consulting task chosen to lie outside the model's competence,
  consultants using AI were 19 percentage points less likely to be
  correct \autocite{DellAcqua2026Navigating}.
\item In a legal examination, untrained lower-ability students adopted
  the model \enquote{but unproductively} \autocite{Chen2026Training}.
\end{itemize}
Outside experiments, one large observational study of about 170,000
programmers found that \enquote{experienced programmers capture nearly
all of these productivity and exploration gains}
\autocite{Daniotti2026Who}. Adoption there was self-selected, and output
was measured as commits. The review cited above sums up the complex end:
\enquote{Across complex tasks, evidence is mixed on whether low or
high-skilled workers benefit more} \autocite{DelRioChanona2025AI}.

\paragraph{Experts on familiar ground; measured versus felt gains.}
Sixteen experienced open-source developers worked on their own mature
projects. They expected AI to speed them up and afterwards estimated a
20\,\% saving, but \enquote{allowing AI actually increases completion
time by 19\%}. They \enquote{slowed down more on issues they are more
familiar with} \autocite{Becker2025Measuring}. In a randomised experiment
with 976 managers, \enquote{experienced managers achieve quality
improvements without confidence inflation, whereas highly educated
managers exhibit confidence gains without corresponding quality
improvements} \autocite{Camuffo2026Beyond}. With models of 2025--2026, a
study of 276 researchers found gains \enquote{across much of the
expertise ladder}. Its point estimates were larger for graduate students
and seniors, but imprecise \autocite{Aparicio2026Generative}.

\section{Discussion and limitations}
\label{who-gains:discussion}

The two strands do not contradict each other once the task is taken into
account. Levelling comes from tasks that are well specified, short and
inside what the model does well: answering support chats, writing
standard texts, completing code, structured consulting tasks. There the
model supplies the missing know-how, and the expert has little left to
gain. Amplification comes from tasks that are open-ended, lie outside
the model's competence, or whose output must be judged: business
decisions, mathematical answers that may be hallucinated, a task chosen
to fool the model. There the gain depends on recognising good output.
The review of the field states the same split \autocite{DelRioChanona2025AI}.
The lawyer and education experiments add a third point. Even where
assisted performance levels, unassisted skill stays with those who had
it \autocite{Autor2026Does,Cruces2026Does}.

The limits of the evidence:
\begin{itemize}
\item Most experiments use short, well-defined tasks, so the levelling
  result is best established exactly where amplification is least
  expected.
\item \enquote{Skill} is measured as tenure, baseline performance,
  education or self-rating, which are not the same thing.
\item Seven cited sources could be read only in their abstracts, and
  four only in an earlier working-paper version beside the published
  abstract: the searches ran without the university network, and two
  publishers answered the later full-text requests with a bot check.
\item DBLP could not be searched. IEEE missed three queries, and two
  answers stopped early.
\item The tier-2 sample was drawn before the repair queries added about
  12,900 records, so it says nothing about those.
\item The screening decisions were made by AI agents and a language
  model and have not yet been checked by a human.
\item Several cited studies are working papers or preprints.
\end{itemize}

\section{Conclusion}
\label{who-gains:conclusion}

On well-defined tasks that the model does well, the less skilled and
less experienced gain more than the experts; on open-ended tasks, tasks
outside the model's competence and tasks whose output must be verified,
those who cannot judge the output gain less or are harmed. The thesis
that AI is in general a skill amplifier is therefore not supported as
stated. The evidence does support a narrower version. What a person gets
out of AI depends on what that person can do without it, but in
opposite directions for the two kinds of task. Durable skill stays with
those who had it in both cases. A single expert's success, such as a
mathematician who can check every step of a proof, illustrates the
second kind of task. It is not evidence about the average worker.

\section{Future work}
\label{who-gains:future}

This search left some work undone. Seven cited studies were read only in
abstract, and four only in an earlier version; they should be read in
full. Among them, the levelling result of \textcite{Cui2026Effects},
the published version of \textcite{Noy2023Experimental} and the
mechanism reported by \textcite{Otis2026Uneven} matter most. The
published version of \textcite{DellAcqua2026Navigating} should be
checked for the 43/17 split. Four relevant works that the semantic search
returned have no indexed abstract and were not obtained:
\begin{itemize}
\item a study of temporary upskilling of knowledge workers;
\item a UK experiment titled as finding \enquote{no great equalizer};
\item a commentary titled \enquote{AI amplifies experts, not expertise};
\item an experiment on why experts miss AI errors.
\end{itemize}
DBLP, the IEEE queries that hit the quota, and the two incomplete
answers should be rerun. Twenty excluded records without abstracts may
bear on the question. A human should review a random sample of the
machine's decisions, and the tier-2 sample should be redrawn over the
records added by the repair queries.

The field lacks experiments on long, open-ended expert work, where
amplification is expected but rarely measured. It also lacks studies
that follow the same workers over months, both with and without the
assistant. If such studies found levelling there too, the lecture's
message would become simply that AI helps beginners most. If they found
amplification, the task distinction drawn here would be confirmed.

\section{Hits that bear on the claim}
\label{who-gains:hits}

% Author's note: generated by
%   scholar sessions export who-gains -f table --bearing-only --lang en \
%     --label who-gains:hits-table --track "Who gains from generative AI" \
%     --theme levelling=... -o litteratursokning/who-gains-full
% raw export beside it (who-gains-full-raw.csv).
% NOTE: caption rewritten by who-gains-screen/fix_caption.py to state what
% was screened and what was only counted; re-run it after a re-export.
% --- scholar LaTeX fragment --------------------------------------------
% Generated by: scholar sessions export -f table
% Session: who-gains
% \input this file from the body of a document whose
% preamble loads:
%   \usepackage{longtable}
%   \usepackage{booktabs}
%   \usepackage{array}
% Text that pdfLaTeX cannot set in T1 (Cyrillic, CJK, ...) is
% wrapped in \scholarnonlatin{...}. The \providecommand lines
% below print the text under XeLaTeX and LuaLaTeX (load a font
% that has the script) and print [non-Latin text] under pdfLaTeX.
% To print it under pdfLaTeX too, define the command before the
% \input, e.g. for Cyrillic with \usepackage[T2A,T1]{fontenc}:
%   \newcommand{\scholarnonlatin}[1]{{\fontencoding{T2A}\selectfont #1}}
% ----------------------------------------------------------------------

\ifdefined\XeTeXversion\providecommand{\scholarnonlatin}[1]{#1}\fi
\ifdefined\directlua\providecommand{\scholarnonlatin}[1]{#1}\fi
\providecommand{\scholarnonlatin}[1]{[non-Latin text]}
\begingroup\footnotesize
\begin{longtable}{@{}%
  >{\raggedright\arraybackslash}p{0.44\textwidth}c%
  >{\raggedright\arraybackslash}p{0.13\textwidth}%
  >{\raggedright\arraybackslash}p{0.26\textwidth}@{}}
\caption{Records bearing on the claim, Who gains from generative-AI assistance (0 cited, 21 supporting and 77 qualifying or contradicting it, out of 24,774 unique hits). What was screened: the first tier (the first 100 hits of every relevance-ordered answer, the works found by a semantic check, and named studies; 1,694 records) was screened in full by AI agents; a seeded random sample of 60 of the 13,340 hits outside the first tier before the repair queries was screened the same way, and none bore on the claim (exact binomial 95\,\% upper bound on the share that could: 6.0\,\%; the sample does not cover the hits that the repair queries added later); a single trial batch of 100 further hits was classified by a language model. The remaining 22,966 unique hits were only counted: they were neither screened nor classified. Not listed: the 1,411 adjacent hits, the 231 false hits, the 22,966 unscreened hits, and 68 hits excluded although they bear on the claim. The last column gives the reason for each inclusion; for rows decided by an agent the agent's confidence is given. Databases searched: OpenAlex (OA; not open access), IEEE Xplore (IEEE; failed for 3 of 9 queries), Web of Science (WoS), Scopus (incomplete for 1 of 9 queries), Semantic Scholar (S2; incomplete for 1 of 9 queries). Searched 2026-11-27.}\label{tab:who-gains:hits-table}\\
\toprule
Title & Year & Database & Reason \\
\midrule
\endfirsthead
\multicolumn{4}{@{}l}{\emph{\tablename~\thetable{} (cont.)}}\\
\toprule
Title & Year & Database & Reason \\
\midrule
\endhead
\midrule \multicolumn{4}{r@{}}{\emph{continued on next page}}\\
\endfoot
\bottomrule
\endlastfoot
AI for Low-Code for AI & 2024 & OA, WoS, Scopus, S2 & supports the claim (0.50) \\
Beyond Black Boxes: Designing and Testing Agentic AI Systems for Strategy & 2026 & OA, Scopus, WoS, S2 & supports the claim (0.65) \\
ChatGPT in Nursing Education: Who Benefits Most? & 2026 & OA, WoS, Scopus & supports the claim (0.80) \\
Does it pay off to use GenAI in the Russian labor market? & 2026 & WoS & supports the claim (0.65) \\
From helping hand to stumbling block: the ChatGPT paradox in competency experiment & 2024 & OA, WoS, Scopus, S2 & supports the claim (0.85) \\
How Much Does AI Impact Development Speed? an Enterprise-Based Randomized Controlled Trial & 2025 & IEEE, Scopus & supports the claim (0.50) \\
Integrating Generative AI into Knowledge Management: The Case of Software Development & 2024 & OA, WoS & supports the claim (0.50) \\
Interaction between software developers and large language model -based tools & 2026 & OA & supports the claim (0.50) \\
Redefining efficiency in post-editing of LLM-generated translation: why expert translators take more time & 2026 & OA, S2 & supports the claim (0.50) \\
Task Expansion with Generative AI: Experimental Evidence on Middle-Skilled, Early-Career Workers & 2026 & S2 & supports the claim (0.50) \\
The Impact of Generative AI Hallucination on Coding Performance: The Moderating Role of Expertise and Group Dynamics & 2025 & OA & supports the claim (0.85) \\
The productivity--learning paradox: a systematic review of programmer experience with generative AI Tools & 2026 & S2, OA & supports the claim (0.60) \\
The role of generative artificial intelligence on labor market: literature review & 2026 & OA, S2 & supports the claim (0.60) \\
The Uneven Impact of Generative AI on Entrepreneurial Performance & 2023 & OA, S2 & supports the claim (0.90) \\
The Uneven Impact of Generative Artificial Intelligence on Entrepreneurial Performance: Evidence from a Field Experiment in Kenya & 2026 & OA, WoS, S2 & supports the claim (0.95) \\
The Widening Gap: The Benefits and Harms of Generative AI for Novice Programmers & 2024 & OA, WoS, Scopus, S2 & supports the claim (0.70) \\
Training for Technology: Adoption and Productive Use of Generative AI in Legal Analysis & 2026 & OA, S2 & supports the claim (0.80) \\
We are not afraid of the wolf! - AI usage attitudes among Hungarian informatics students & 2025 & WoS & supports the claim (0.60) \\
What Skills Do You Need When Developing Software Using ChatGPT? (Discussion Paper) & 2023 & S2, OA, Scopus & supports the claim (0.50) \\
Who Benefits More? The Role of AI-Human Interaction in Creativity Augmentation & 2024 & OA & supports the claim (0.80) \\
Who is using AI to code? Global diffusion and impact of generative AI & 2025 & S2, OA, WoS, Scopus & supports the claim (0.90) \\
\scholarnonlatin{Исследование отечественного и глобального рынка консалтинговых услуг: эволюция и современное состояние} (doi:10.18334/epp.16.6.124921) & 2026 & OA & qualifies or contradicts the claim (0.90) \\
`` ChatGPT can make mistakes. Check important info.'' Epistemic beliefs and metacognitive accuracy in students' integration of ChatGPT content into academic writing & 2025 & OA, WoS, Scopus, S2 & qualifies or contradicts the claim (0.50) \\
A Mixed-Method Empirical Study of LLM Assistance in Software Engineering Workflows & 2026 & OA, Scopus, S2 & qualifies or contradicts the claim (0.70) \\
A Multiagent Large Language Model Framework for Emergency Treatment Recommendation in Acute Ischemic Stroke: Development and Validation Study & 2026 & OA, Scopus, WoS & qualifies or contradicts the claim (0.85) \\
A pilot study on the impact of large language model assistance on the evaluation of complex medical living kidney donor candidates & 2026 & OA, Scopus, S2 & qualifies or contradicts the claim (0.70) \\
A Systematic Evaluation of Code-generating Chatbots for Use in Undergraduate Computer Science Education & 2024 & IEEE, WoS, Scopus, OA, S2 & qualifies or contradicts the claim (0.55) \\
AI and Freelancers: Has the Inflection Point Arrived? & 2024 & Scopus & qualifies or contradicts the claim (0.55) \\
AI and jobs. A review of theory, estimates, and evidence & 2025 & OA, S2 & qualifies or contradicts the claim (0.80) \\
AI as a Procedural Equalizer: Performance Comparison in Programming-Based Engineering Coursework Following the Emergence of Generative AI & 2026 & WoS & qualifies or contradicts the claim (0.50) \\
AI at the bedside: Randomised controlled trial of ChatGPT's impact on student performance in real-patient clinical exams & 2026 & WoS & qualifies or contradicts the claim (0.65) \\
AI TOOLS IN THE WORKPLACE: PRODUCTIVITY GAINS, DESKILLING RISKS AND THE ROLE OF CONTINUOUS EMPLOYEE EDUCATION & 2026 & OA, S2 & qualifies or contradicts the claim (0.65) \\
Artificial Intelligence and Jobs: Has the Inflection Point Arrived? Evidence from an Online Labor Platform & 2026 & OA, WoS, Scopus, S2 & qualifies or contradicts the claim (0.55) \\
Automation as an Equalizer: How EasytoUse Technologies Narrow Skill Gaps Between Low and HighSkilled Workers & 2026 & OA, WoS, Scopus, S2 & qualifies or contradicts the claim (0.60) \\
Beyond Code: The Multidimensional Impacts of Large Language Models in Software Development & 2025 & OA, S2 & qualifies or contradicts the claim (0.65) (no-abstract) \\
Bridging the Expertise Gap: The Role of Generative AI in Supporting Project Planning Tasks for Novices and Professionals & 2025 & OA, WoS, S2, Scopus & qualifies or contradicts the claim (0.90) \\
Bridging the Skill Gap in Clinical CBCT Interpretation with CBCTRepD & 2026 & OA, S2 & qualifies or contradicts the claim (0.60) \\
Can a troubleshooting AI assistant improve task performance in industrial contexts? & 2025 & OA, WoS, Scopus & qualifies or contradicts the claim (0.85) \\
Can ChatGPT Kill User-Generated Q\&A Platforms? & 2026 & WoS & qualifies or contradicts the claim (0.50) \\
Copilot Impact Studies: Measuring Productivity, Trust, and Skill Evolution in Enterprise Developer Teams & 2025 & OA, S2 & qualifies or contradicts the claim (0.50) \\
Does AI Assistance Enhance or Erode Expertise? Evidence from a Three-Month Field Experiment in Patent Drafting & 2026 & OA & qualifies or contradicts the claim (0.60) \\
Does AI Democratize Commercial Content Creation? Changes to the Strategic Communication Profession & 2025 & OA, Scopus & qualifies or contradicts the claim (0.80) \\
Does Generative AI Narrow Education-Based Productivity Gaps? Evidence from a Randomized Experiment & 2026 & OA, S2 & qualifies or contradicts the claim (0.90) \\
Effectiveness of LLM-based AI Assistance for small business productivity & 2025 & OA & qualifies or contradicts the claim (0.85) \\
Engaging Customers with AI in Online Chats: Evidence from a Randomized Field Experiment & 2025 & OA, Scopus, S2 & qualifies or contradicts the claim (0.90) \\
Enhancing sonologist examination performance with large language models: an analytical study of ChatGPT-4 and Claude 3 & 2025 & OA, WoS, Scopus, S2 & qualifies or contradicts the claim (0.50) \\
Evaluation of large language model-driven AutoML in data and model management from human-centered perspective & 2025 & OA, WoS, Scopus, S2 & qualifies or contradicts the claim (0.50) \\
Experimental evidence on the productivity effects of generative artificial intelligence & 2023 & S2 & qualifies or contradicts the claim (0.95) \\
Experimenting with Generative AI: Does ChatGPT Really Increase Everyone's Productivity? & 2025 & OA, S2 & qualifies or contradicts the claim (0.75) \\
Exploring Large Language Model AI tools in Construction Project Risk Assessment: Chat GPT Limitations in Risk Identification, Mitigation Strategies, and User Experience & 2025 & OA, WoS, Scopus & qualifies or contradicts the claim (0.55) \\
Generating creativity through ChatGPT: an empirical investigation in open innovation platforms & 2025 & OA, Scopus, WoS, S2 & qualifies or contradicts the claim (0.75) \\
Generative AI and labour productivity: A quasi experiment on coding & 2026 & OA, WoS, Scopus & qualifies or contradicts the claim (0.95) \\
Generative AI and Technical Debt in Software Development: Productivity Gains versus Code Quality and Security Risk & 2026 & OA, S2 & qualifies or contradicts the claim (0.55) \\
Generative AI and the Productivity Divide: Human-AI Complementarities in Education & 2026 & OA, S2 & qualifies or contradicts the claim (0.65) \\
Generative AI at Work & 2025 & OA, WoS, Scopus, S2 & qualifies or contradicts the claim (0.97) \\
Generative AI Boosts Productivity Across the Expertise Ladder & 2026 & OA, S2 & qualifies or contradicts the claim (0.65) \\
Generative AI in Action: Field Experimental Evidence from Alibaba's Customer Service Operations & 2026 & OA, S2 & qualifies or contradicts the claim (0.95) \\
Generative AI in public administration: A quasi-experimental analysis of bureaucratic productivity & 2026 & OA, WoS, Scopus & qualifies or contradicts the claim (0.85) \\
Generative AI in the Workplace: A Systematic Review of Productivity Effects, Employment Perceptions, and Job Insecurity & 2026 & OA, S2 & qualifies or contradicts the claim (0.50) \\
Generative AI-assisted workflows in architectural conceptual design: Performance, creative self-efficacy, and cognitive load & 2026 & OA, Scopus, WoS, S2 & qualifies or contradicts the claim (0.60) \\
Generative artificial intelligence produces differential academic outcomes across socioeconomic, linguistic, and disability contexts & 2026 & OA & qualifies or contradicts the claim (0.60) \\
How Artificial Intelligence Technology Affects Productivity & 2025 & OA, Scopus & qualifies or contradicts the claim (0.55) \\
How do copilot suggestions impact developers' frustration and productivity? & 2026 & OA, WoS, S2 & qualifies or contradicts the claim (0.85) \\
How far are AI-powered programming assistants from meeting developers' needs? & 2024 & OA, S2 & qualifies or contradicts the claim (0.70) \\
How Generative AI Transforms Questioning Behavior on Q\&A Platforms: Evidence from a Natural Experiment with Pilot Usage of ChatGPT & 2026 & OA, S2 & qualifies or contradicts the claim (0.55) \\
How generative artificial intelligence (AI) tools can enhance software development? & 2026 & OA & qualifies or contradicts the claim (0.50) \\
HUMAN--AI COLLABORATION IN KNOWLEDGE WORK: PRODUCTIVITY, ERRORS, AND ETHICAL RISK & 2025 & OA, S2 & qualifies or contradicts the claim (0.70) \\
Is AI Productivity Mediated by User Skill? Prompting Skill, Subject Knowledge, and Human Capital in AI-Assisted Knowledge Production & 2026 & OA & qualifies or contradicts the claim (0.55) \\
Large Language Model in Creative Work: The Role of Collaboration Modality and User Expertise & 2024 & Scopus, WoS, OA & qualifies or contradicts the claim (0.90) \\
Large Language Models in Academia: Boosting Productivity but Reinforcing Inequality & 2025 & OA & qualifies or contradicts the claim (0.90) \\
Lumos: AI-driven prompt optimisation tool for assisting conceptual design & 2024 & Scopus, WoS & qualifies or contradicts the claim (0.60) \\
Measuring the Impact of Early-2025 AI on Experienced Open-Source Developer Productivity & 2025 & OA, S2 & qualifies or contradicts the claim (0.75) \\
Navigating the Jagged Technological Frontier: Field Experimental Evidence of the Effects of Artificial Intelligence on Knowledge Worker Productivity and Quality & 2026 & WoS & qualifies or contradicts the claim (0.85) \\
Not Faster, but Better: Generative AI and Teacher Skill in Peru & 2026 & OA & qualifies or contradicts the claim (0.65) \\
Personalized Coding Assistants Adapting Large Language Models to Individual Developer Styles & 2023 & OA, S2 & qualifies or contradicts the claim (0.50) \\
Promotion or substitution? The effect of AI assistants on physician performance in online healthcare platforms & 2026 & Scopus, WoS, OA, S2 & qualifies or contradicts the claim (0.60) \\
PsychBench: A comprehensive and professional benchmark for evaluating the performance of LLM-assisted psychiatric clinical practice & 2025 & OA, S2 & qualifies or contradicts the claim (0.80) \\
Shaping Human-AI Collaboration: Varied Scaffolding Levels in Co-writing with Language Models & 2024 & WoS & qualifies or contradicts the claim (0.75) \\
Short-Term Efficiency, Long-Term Cognitive Cost? A Systematic Literature Review of Large Language Model Use in Professional Work & 2026 & OA & qualifies or contradicts the claim (0.90) \\
Structured AI Integration for Equitable STEM Writing: A Pilot Study of the DISCO Framework & 2026 & OA, Scopus & qualifies or contradicts the claim (0.60) \\
The ABC's of Who Benefits from Working with AI: Ability, Beliefs, and Calibration & 2024 & OA, WoS, Scopus & qualifies or contradicts the claim (0.95) \\
THE ASYMMETRIC IMPACT OF GENERATIVE ARTIFICIAL INTELLIGENCE ON THE RETURNS TO COGNITIVE LABOR: A QUASI-EXPERIMENTAL STUDY OF HIGH-SKILLED FREELANCE MARKETS & 2026 & OA & qualifies or contradicts the claim (0.60) \\
The Double-Edged Roles of Generative AI in the Creative Process: Experiments on Design Work & 2025 & WoS, OA, S2 & qualifies or contradicts the claim (0.90) \\
The Economic Impact of ChatGPT as a Digital Assistant on Workplace Productivity & 2026 & OA, Scopus & qualifies or contradicts the claim (0.70) \\
The Effect of Generative Artificial Intelligence on Software Engineering Productivity in the Software Industry & 2026 & IEEE, OA, WoS, Scopus, S2 & qualifies or contradicts the claim (0.60) \\
The Effects of Generative AI on High-Skilled Work: Evidence from Three Field Experiments with Software Developers & 2026 & OA, WoS, S2 & qualifies or contradicts the claim (0.95) \\
The Impact of Generative AI on Knowledge Work and Student Learning: A Systematic Literature Review & 2026 & OA & qualifies or contradicts the claim (0.60) \\
The Invisible Gap: How AI Productivity Masks Eroding Expertise in Knowledge Work -- and what to do about it & 2026 & OA & qualifies or contradicts the claim (0.50) \\
The Role of Artificial Intelligence and Professional Expertise in Adapted Physical Activity Prescription for Orthopedic Rehabilitation & 2026 & OA, WoS, Scopus, S2 & qualifies or contradicts the claim (0.50) \\
The Short-Term Effects of Generative Artificial Intelligence on Employment & 2024 & OA & qualifies or contradicts the claim (0.50) \\
The Short-Term Effects of Generative Artificial Intelligence on Employment: Evidence from an Online Labor Market & 2024 & S2, OA, WoS, Scopus & qualifies or contradicts the claim (0.50) \\
Trusta: Reasoning about assurance cases with formal methods and large language models & 2025 & OA, WoS, Scopus & qualifies or contradicts the claim (0.55) \\
Tutor CoPilot: A Human-AI Approach for Scaling Real-Time Expertise & 2024 & OA, S2 & qualifies or contradicts the claim (0.90) \\
UGenLah: A Human-in-the-loop Agentic Paradigm for Game Development & 2026 & WoS, Scopus, S2 & qualifies or contradicts the claim (0.50) \\
Uneven Gains of Generative Artificial Intelligence on Creativity Enhancement: A Systematic Review and Meta-Analysis & 2025 & OA, S2 & qualifies or contradicts the claim (0.85) \\
Use of Generative Artificial Intelligence & 2025 & OA & qualifies or contradicts the claim (0.95) \\
When Young Scholars Cooperate with LLMs in Academic Tasks: The Influence of Individual Differences and Task Complexities & 2025 & WoS & qualifies or contradicts the claim (0.50) \\
Who listens to ChatGPT and when should they? A two-study examination of AI-assisted decision making & 2026 & OA, WoS & qualifies or contradicts the claim (0.55) \\
\end{longtable}
\endgroup

